\documentclass{article}
\usepackage[T1]{fontenc}
\usepackage{lmodern}
\usepackage{graphicx}
\usepackage{amsmath,amssymb}
\usepackage[a4paper,margin=2cm]{geometry}
\usepackage[absolute,overlay]{textpos}
\setlength{\TPHorizModule}{1mm}
\setlength{\TPVertModule}{1mm}
\usepackage[indonesian]{babel}
\usepackage{hyperref}
\setlength{\emergencystretch}{2em}
% unit: Halaman 77

\begin{document}

% === Halaman 77 (python/native) ===

Kelebihan Trivium

1. Trivium cocok untuk lightweight cryptography, karena state relatif 

kecil (hanya 288 bit)

2. Hardware efficiency. Sederhana jika diimplementasikan secara 

hardware, karena struktunta sederhana (terutama terdiri dari flip-

flop/register, gerbang AND dan OR)

77

\end{document}
